\documentclass[12pt]{article}
\usepackage[T1]{fontenc}
\usepackage[utf8]{inputenc}
\usepackage{lmodern}
\usepackage{textcomp}
\usepackage{enumitem}
\usepackage{geometry}
\usepackage{graphicx}
\usepackage{amsmath, amssymb}
\geometry{margin=1in}
\begin{document}
\begin{center}
Geography - Undergraduate Level Assessment 8 \
Total Marks: 40 \
Answer all questions.
\end{center}

\section*{Multiple Choice (10 marks)}
\begin{enumerate}[label=\arabic*.]
\item As of 2017, the share of GDP spent on the military by the United States is about
\begin{enumerate}[label=(\alph*)]
    \item 1\%
    \item 3\%
    \item 6\%
    \item 10\%
\end{enumerate}
\item As of 2017, the share of deaths in Greenland by suicide is about
\begin{enumerate}[label=(\alph*)]
    \item 0.90\%
    \item 1.80\%
    \item 3.60\%
    \item 7.20\%
\end{enumerate}
\item What is the percentage of children aged 13-15 in Japan who reported being bullied at least once in the past couple of months as of 2015?
\begin{enumerate}[label=(\alph*)]
    \item 6\%
    \item 16\%
    \item 36\%
    \item 66\%
\end{enumerate}
\item As of 2014, about what percentage of primary aged children in Sub-Saharan Africa were out- of-school?
\begin{enumerate}[label=(\alph*)]
    \item 10\%
    \item 20\%
    \item 35\%
    \item 50\%
\end{enumerate}
\item In 2017, about how many people died from terrorism globally?
\begin{enumerate}[label=(\alph*)]
    \item 260
    \item 2,600
    \item 26,000
    \item 260,000
\end{enumerate}
\end{enumerate}

\section*{True / False (10 marks)}
\textit{Answer True or False}
\begin{enumerate}[label=\arabic*.]
\setcounter{enumi}{5}
\item True or False: It was an area of political and geographical control imposed by the Viking incursions which was formally agreed to by the Danish warlord, Guthrum and west-Saxon king, a different answer in 886 AD.
\item True or False: Houston has sports teams for every major professional league except the National Hockey League (NHL). The Houston Astros are a Major League Baseball (MLB) expansion team formed in 1962 (known as the "a different answer" until 1965) that made one World Series appearance in 2005.
\item True or False: These clashes, in part, contributed to the Panic of 1873 which delayed construction of the railroad into Montana. Surveys in 1874, 1875 and 1876 helped spark the Great Sioux War of 1876. The transcontinental NPR was completed on September 8, 2005, at Gold Creek.
\item True or False: Southampton Solent University has 18539 students and its strengths are in the training, design, consultancy, research and other services undertaken for business and industry.
\item True or False: The New York City Fire Department is the largest municipal fire department in the United States and the second largest in the world after the Tokyo Fire Department. The FDNY employs approximately 11,080 uniformed firefighters and over 3,300 uniformed EMTs and paramedics. The FDNY's motto is a different answer.
\end{enumerate}

\section*{Long Form Response (20 marks)}
\textit{Answer each question with a clear and complete explanation.}
\begin{enumerate}[label=\arabic*.]
\setcounter{enumi}{10}
\item Read the following passage and answer the question: At no more than 200 kilometres (120 mi) north to south and 130 kilometres (81 mi) east to west, Swaziland is one of the smallest countries in Africa. Despite its size, however, its climate and topography is diverse, ranging from a cool and mountainous highveld to a hot and dry lowveld. The population is primarily ethnic Swazis whose language is siSwati. They established their kingdom in the mid-18 th century under the leadership of Ngwane III; the present boundaries were drawn up in 1881. After the Anglo-Boer War, Swaziland was a British protectorate from 1903 until 1967. It regained its independence on 6 September 1968. Question: What is the primary language spoken by the people in Swaziland?
\item Read the following passage and answer the question: Napoleon was born on 15 August 1769, to Carlo Maria di Buonaparte and Maria Letizia Ramolino, in his family's ancestral home Casa Buonaparte in Ajaccio, the capital of the island of Corsica. He was their fourth child and third son. This was a year after the island was transferred to France by the Republic of Genoa. He was christened Napoleone di Buonaparte, probably named after an uncle (an older brother who did not survive infancy was the first of the sons to be called Napoleone). In his 20 s, he adopted the more French-sounding Napoléon Bonaparte.[note 2] Question: Napoleon was most likely named after what family relation?
\end{enumerate}

\vfill
\noindent\textit{End of Paper}
\end{document}